\documentclass[11pt,a4paper,twoside]{tesis}
% SI NO PENSAS IMPRIMIRLO EN FORMATO LIBRO PODES USAR
%\documentclass[11pt,a4paper]{tesis}

\usepackage{graphicx}
\usepackage[utf8]{inputenc}
\usepackage[spanish]{babel}
\usepackage[left=3cm,right=3cm,bottom=3.5cm,top=3.5cm]{geometry}

\begin{document}

%%%% CARATULA

\def\autor{Adrian Norberto Marino}
\def\tituloTesis{Sistemas de recomendación colaborativos e híbridos}
\def\runtitulo{Resumen}
\def\runtitle{Sistemas de recomendación colaborativos e híbridos}
%\def\director{Obi-Wan Kenobi}
%\def\codirector{Master Yoda}
\def\lugar{Buenos Aires, 2022}
\newcommand{\HRule}{\rule{\linewidth}{0.2mm}}
%
\thispagestyle{empty}

\begin{center}\leavevmode

\vspace{-2cm}

\begin{tabular}{l}
\includegraphics[width=3.6cm]{logouba.png}
\end{tabular}


{\large \sc Universidad de Buenos Aires\\Facultad de Ciencias Exactas y Naturales \\ Facultad de Ingeniería}

\vspace{6.0cm}

\vspace{3.0cm}
%{
%\Large \color{red}
%\begin{tabular}{|p{2cm}cp{2cm}|}
%\hline
%& Pre-Final Version: \today &\\
%\hline
%\end{tabular}
%}
%\vspace{2.5cm}

\begin{huge}
\textbf{\tituloTesis}
\end{huge}

\vspace{2cm}

{\large Trabajo Final de Especialización en Explotación de Datos y Descubrimiento del Conocimiento}

\vspace{2cm}

{\Large \autor}

\end{center}

\vfill

{\large

%{Director: \director}

\vspace{.2cm}

%{Codirector: \codirector}

\vspace{.2cm}

\lugar
}

\newpage\thispagestyle{empty}

%
%%%% ABSTRACTS, AGRADECIMIENTOS Y DEDICATORIA
\frontmatter
\pagestyle{empty}
%\begin{center}
%\large \bf \runtitulo
%\end{center}
%\vspace{1cm}
\chapter*{\runtitulo}

\noindent Acá va un resumen del trabajo. Con el resumen se deberia poder tener una idea del trabajo en su totalidad. Desde los obetivos y los datos hasta los resultados y conclusiones. Se escibe al final. (aprox. 200 palabras).

\bigskip

\noindent\textbf{Palabras claves:} representativas del trabajo, los métodos, los datos (no menos de 5).

%\cleardoublepage
%%\begin{center}
%\large \bf \runtitle
%\end{center}
%\vspace{1cm}
\chapter*{\runtitle}

\noindent In a galaxy far, far away, a psychopathic emperor and his most trusted servant -- a former Jedi Knight known as Darth Vader -- are ruling a universe with fear. They have built a horrifying weapon known as the Death Star, a giant battle station capable of annihilating a world in less than a second. When the Death Star's master plans are captured by the fledgling Rebel Alliance, Vader starts a pursuit of the ship carrying them. A young dissident Senator, Leia Organa, is aboard the ship \& puts the plans into a maintenance robot named R2-D2. Although she is captured, the Death Star plans cannot be found, as R2 \& his companion, a tall robot named C-3PO, have escaped to the desert world of Tatooine below. Through a series of mishaps, the robots end up in the hands of a farm boy named Luke Skywalker, who lives with his Uncle Owen \& Aunt Beru. Owen \& Beru are viciously murdered by the Empire's stormtroopers who are trying to recover the plans, and Luke \& the robots meet with former Jedi Knight Obi-Wan Kenobi to try to return the plans to Leia Organa's home, Alderaan. After contracting a pilot named Han Solo \& his Wookiee companion Chewbacca, they escape an Imperial blockade. But when they reach Alderaan's coordinates, they find it destroyed - by the Death Star. They soon find themselves caught in a tractor beam \& pulled into the Death Star. Although they rescue Leia Organa from the Death Star after a series of narrow escapes, Kenobi becomes one with the Force after being killed by his former pupil - Darth Vader. They reach the Alliance's base on Yavin's fourth moon, but the Imperials are in hot pursuit with the Death Star, and plan to annihilate the Rebel base. The Rebels must quickly find a way to eliminate the Death Star before it destroys them as it did Alderaan (aprox. 200 palabras).

\bigskip

\noindent\textbf{Keywords:} War, Rebellion, Wookie, Jedi, The Force, Empire (no menos de 5). % OPCIONAL: comentar si no se quiere

%\cleardoublepage
%\chapter*{Agradecimientos}

\noindent  % OPCIONAL: comentar si no se quiere

%\cleardoublepage
%\hfill \textit{A mi persona favorita.}
  % OPCIONAL: comentar si no se quiere

%\cleardoublepage
\tableofcontents

\mainmatter
\pagestyle{headings}

%%%% ACA VA EL CONTENIDO DE LA TESIS

\chapter{Introducción}

Los sistemas de recomendación tienen por objetivo acercar a sus usuarios información, productos, contenido (Textos,  audio, videos), etc.. relevantes a sus preferencias o necesidades, permitiendo a estos encontrar con mayor facilidad aquello que buscan.  Formalizando esta definición podemos decir que:  Los sistemas de recomendación apuntan a ayudar a un usuario o grupo de usuarios a seleccionar items de forma personalizada dado un de conjunto de items de gran extensión o un gran espacio de búsqueda.


Este objetivo puede cambiar según el contexto de cada negocio. Para un e-commerce de delivery de comidas, el objetivo es acercar a los usuarios el tipo de comida que quieren probar en ese mismo momento, a un precio que puedan pagar con tiempo de entrega aceptable. Para un e-commerce de venta de productos, se busca acercar al usuario aquellos productos que este necesitando en ese mismo momento, los cuales tienen un precio que el mismo puede pagar y por otro lado, asegurar una experiencia satisfactoria con el vendedor. En el negocio de visualización de contenido (Ya sea audio o video), el objetivo es acercar al usuario contenido a fin a sus gustos para mejorar su experiencia en la plataforma y así aumentar el engagement de sus usuarios.


Por otro lado, el objetivo de fondo siempre es el mismo, mejorar la conversión. Con esto nos referimos a aumentar el volumen de ventas para un e-commerce de venta de productos, la cantidad de deliveries mensuales, la cantidad de impresiones de publicidad  en aplicaciones de visualización de contenido, aumentar el tiempo de permanecía en las plataformas de streaming de audio o video, etc.. Podemos encontrar muchos ejemplos distintos donde el objetivo común es mejorar la conversión y engagement de los usuarios.

Desde un punto de vista mas técnico, los sistemas de recomendación se utilizan para predecir el grado de preferencia de un usuario con respecto a un item. 
En general, se puede lograr aplicando de un algoritmo de optimización, el cual minimiza la diferencia entre el grado de preferencia esperado versus real. Otros enfoques hace uso de medidas de distancia para establecer este grado de preferencia.

\section{Tipos de sistemas de recomendación}


\begin{itemize}

\item Popularidad: Este tipo de recomendador toma alguna caracteristica de popularidad de los items en cuestión, como pueden ser: cantidad de vistas, cantidad de compras, cantidad de reviews positivos, etc.. Luego se ranquea el top N de items mas populares. Si bien este tipo de recomendacion tiene buenos resutaldos para usaurios nuevos, de los cuales no se conocen sus gustos, carecen de personalización, ya que las recomendaciones no tienen en cuenta las preferencias de cada usuario.

\item Basados en contenido: A diferencia de los recomendadores basados en filtros colaborativos, este tipo necesita un trabajo trabajo previo de ingenieria de features sobre los items, donde se busca definir que features son los mas significativos para la tarea y cuales el grado de adecuacion de cada items a estaod features. Por otro lado, es necesario registrar las interacciones de los usuario. luego con estas interacciones se puede definir el grado de preferencia de los usuario sobre sobre los features de los items. Con esta informacion se puede encontrar tanto items como usaurios similares y realizar recomendaciones del tipo:

\begin{itemize}

\item Si el usuario A tiene preferencia por el item X , seguramente le gustara el item B por ser muy cercano al item A.
\item Si el usuario A y el usuario B son muy cercanos en cuando al grado preferencias, es posible recomendar items consumidos el usaurio B al usaurio A y vise versa.

\end{itemize}


y por otro lado el grado de preferencia de cada usuario por la

\item Basados en filtrado contenido

\item Híbridos

\end{itemize}	


\begin{itemize}
\item Considerando los datos a utilizar: Estos dependen de que datos utilicemos para construirlos.	
	\begin{itemize}
	\item Basados en reglas: Son modelos que halla un árbol de decisión en base a los datos. A estos arboles podemos verlos como las reglas a utilizar para realizar una recomendación. Como principal ventaja, tienen el mayor grado de explicabilidad posible,  ya que a simple vista se puede apreciar las reglas que el modelo ha encontrado en base a los datos. Como ejemplo, una posible regla hallada podría ser: Si el usuario 1 compro el item A, entonces comprara el item B.
	\item Basados en contenido:  Estos modelos utilizan los atributos o features de los items a recomendar. Por ejemplo, si estamos recomendando películas, los atributos a elegir podrían ser: título, género, director, actores de reparto, fecha de estreno, idioma original, etc... Estos modelos  en general codifican estos features usando distintas técnicas (One-hot encoding, Embeddings, TF-IDF, etc..), en algunas de esta tecnicas se busca bajar la dimensionalidad de los mismos sin perder información, en otras solo codificar la informacion a un formato que entiendan los modelos. 
De esta forma, se generan vectores que representan a cada item y finalmente usando alguna medida de distancia se establece el grado de cercanía entre estos vectores. Esto permite resuelver  problemas del tipo: dato un item, cuales son los items mas parecidos o cercanos.
	\item Basados en filtros colaborativos
\end{itemize}	


\item Considerando el modelo: Se refirere a la forma en la que se realiza la predicción.	
	\begin{itemize}
	\item Basados en memoria: Este tipo de modelos mantiene sus datos en memoria y clasificacion recorrer todos los datos(full scan) cada vez que necesita realizar un inferencia o predicción. Un ejemplo es el algoritmo de k vecinos cercanos (KNN), el cual mantiene una matriz rala de distancias en memoria, la cual recorre completamente para compara las distancias entre filas o columnas usando alguna media de distancia como puede ser la Distancia Coseno, Distancia Coseno Ajustada, Distancia de Manhattan, etc...
	\item Basados en modelos:  Estos están constituidos por un modelo basados en los datos el cual se reutiliza cada vez que se desea realizar un inferencia o predicción (Por ejemplo: Single Value Descomposition)
	\end{itemize}
\end{itemize}

\section{Descripción del problema y motivación ($>$ 1/2 carilla)}

Ejemplo de ecuación:
\begin{equation}
L = - \sum_{p \in P} \sum_{y \in Y} w_{py} \sum_{t \in T_{py}} y_t \log \hat y_t + (1-y_t) \log (1-\hat y_t)
\label{eq:loss}
\end{equation}

\section{Trabajos previos ($>$ 1 carilla)}
En la imagen \ref{fig:label} hay un ejemplo de como poner una figura:
\begin{figure}[h]
\centering
\includegraphics[width=0.5\columnwidth]{logouba.png}
\caption{Acá va una explicación de la figura. Recuerden usar colores solo si son explicativos, no olvidar nombrar correctamente los ejes y poner una explicación en el caption.}
\label{fig:label}
\end{figure}

\section{Objetivos ($>$ 1/2 carilla)}

\chapter{Materiales y Métodos}

\section{Datos ($>$ 1 carilla)} 
Descripción de los datos / obligatorio poner la fuente.

\section{Análisis exploratorio}

\chapter{Métodos ($>$ 3 carillas)}
Describir los métodos utilizados (a utilizar). 
Es bligatorio agregar un parrafo especificando a que materías se relacionan.

\section{Método I }
\section{Método II}
\section{Métricas}

\chapter{Experimentos (si los hay)}
Pueden ser comparaciones entre los métodos utilizados, por ejemplo. 

En la \ref{table:tab} se ve un ejemplo de tabla:
\begin{table}[h!]
\centering
\footnotesize
\begin{tabular}{lrrrrrr}
\hline
Phone &  Total & \% Errors &F1 &  EER \\
\hline
   EY &    441 &  13.83 & 0.80 & 0.12 \\
   JH &    178 &  39.89 & 0.81 & 0.18 \\
   AY &   1040 &   5.96 & 0.45 & 0.21 \\
    R &   1298 &  18.34 & 0.53 & 0.27 \\
\hline
\end{tabular}
\caption{Describir la tabla como para que alguien que agarra el trabajo no dependa del cuerpo del texto para entenderla}
\label{table:tab}
\end{table}

\chapter{Resultados ($>$ 3 carillas)}
Pueden ser preliminares. Es bligatorio tener algo hecho.


\chapter{Conclusiones ( $>$ 1 carilla)} 
Pueden ser preliminares. 
Si no hay mucho hecho se pueden discutir las dificultades a futuro.


%%%% BIBLIOGRAFIA
\backmatter
%\bibliography{tesis}

\end{document}
